% ==============================================================================
% 前言
% ==============================================================================

\chapter*{前言}
\addcontentsline{toc}{chapter}{前言}

\section*{编制背景}

随着汽车产业的快速发展，特别是新能源汽车、智能网联汽车的快速普及，汽车外部照明与光信号装置的技术形态正在经历深刻变革。从传统的卤素前照灯到 LED 前照灯，再到矩阵式自适应远光灯（ADB）和自适应前照灯系统（AFS），汽车照明技术持续演进，对安装规定的及时更新提出了迫切要求。

在此背景下，中国国家标准 GB 4785《汽车及挂车外部照明和光信号装置的安装规定》于 2019 年发布实施以来，已运行六年有余。为了与同期 UNECE（联合国欧洲经济委员会）第 48 号法规（UN R48）的最新版本保持一致，应对智能照明新技术的发展需求，国家标准化主管部门启动了 GB 4785 标准的修订工作，并于 2026 年 8 月形成了 V4.3.5 版本的征求意见稿。

本报告应客户要求，针对以下三个版本开展系统的横向对比研究：

\begin{itemize}[leftmargin=2em]
    \item \textbf{GB 4785-2019}：中国现行版本，等效采用 UN R48 系列修订并结合国情作出修改；
    \item \textbf{GB 4785 草案 V4.3.5}：征求意见阶段的最新版本，标志着中国标准即将进入新一轮升级；
    \item \textbf{UN R48-09}：UNECE 在 2025 年 9 月 26 日正式生效的最新修订系列，代表国际最先进的技术要求。
\end{itemize}

\section*{编制目的}

本报告旨在通过系统化的对比分析，回答以下关键问题：

\begin{enumerate}[leftmargin=2em]
    \item 三个版本在适用范围、章节结构、技术要求上有哪些差异？
    \item 草案版相对现行版，新增、删除、修改了哪些关键条款？
    \item 中国标准（现行版和草案版）与欧洲标准（UN R48-09）的接轨程度如何？
    \item 草案版向 UN R48-09 的对齐方向是什么？还存在哪些技术差距？
    \item 这些变化对整车厂、灯具供应商、检测认证机构的影响是什么？
\end{enumerate}

\section*{适用范围}

本报告适用于以下读者：
\begin{itemize}[leftmargin=2em]
    \item 整车制造企业的法规合规团队和技术研发人员；
    \item 汽车灯具供应商的产品规划和研发团队；
    \item 检测认证机构的工程师和审查员；
    \item 行业研究机构和高校研究人员；
    \item 汽车行业政策制定和标准研究的相关人员。
\end{itemize}

\section*{编制原则}

本报告遵循以下编制原则：

\begin{description}[leftmargin=2em, style=nextline]
    \item[系统性]  对三版本进行全章节、全条款的系统对比，不留遗漏；
    \item[准确性]  所有数据、术语、编号均严格遵循官方法规文本；
    \item[实用性]  对比结论以工程应用为导向，提供可操作的合规指引；
    \item[客观性]  对比分析基于法规文本事实，避免主观臆断；
    \item[前瞻性]  在结论中关注技术发展趋势和未来修订方向。
\end{description}

\section*{特别说明}

本报告中所引用的具体条款编号（如 GB 4785-2019 第 5.2.1 条）为法规文本中的原始编号。在实际对比中，章节编号可能在不同版本之间存在调整。本报告在每个对比表格中均明确标注所参考的版本，以确保可追溯性。

报告中标注符号约定：
\begin{itemize}[leftmargin=2em]
    \item \textcolor{themegreen}{\textbf{[新增]}}——在新版本中新增加的、此前版本中没有的条款；
    \item \textcolor{themered}{\textbf{[删除]}}——在新版本中被删除的、此前版本中存在的条款；
    \item \textcolor{themeorange}{\textbf{[修改]}}——条款存在但内容发生了实质性调整；
    \item \textcolor{themeblue}{\textbf{[保留]}}——条款基本维持原有内容。
\end{itemize}

\vspace{1cm}
\begin{flushright}
\begin{minipage}{0.4\textwidth}
\textit{汽车法规对比研究团队}\\
\textit{2026 年 9 月}
\end{minipage}
\end{flushright}

\clearpage

% ==============================================================================
% 摘要
% ==============================================================================

\chapter*{摘要}
\addcontentsline{toc}{chapter}{摘要}

\section*{研究概述}

本报告对 GB 4785-2019（现行版）、GB 4785 草案 V4.3.5（草案版）和 UN R48-09（欧版）三个版本的汽车外部照明与光信号装置安装规定进行系统的横向对比研究。重点关注三者之间的技术差异、演进方向和合规影响。

\section*{核心发现}

通过逐章节、逐条款的对比分析，本报告得出以下核心发现：

\textbf{一、整体架构层面：} 三版本在总体框架上保持高度相似性。草案版基本沿用了现行版的章节结构，但更加紧密地与 UN R48-09 对齐，新增了部分章节以反映新技术的发展需求。

\textbf{二、关键技术变化：}
\begin{itemize}[leftmargin=2em]
    \item 新增自适应远光灯（ADB）系统的安装与控制要求；
    \item 强化自适应前照灯系统（AFS）的等级转换与故障指示要求；
    \item 调整昼间行驶灯（DRL）的安装数量与光度要求；
    \item 新增紧急制动信号（ESS）的强制安装要求；
    \item 完善对 LED 光源及矩阵式 LED 灯具的特殊要求；
    \item 修订部分色度坐标以匹配最新国际标准；
    \item 强化电磁兼容性（EMC）与电气安全要求。
\end{itemize}

\textbf{三、与 UN R48-09 的对齐度：} 草案版相对现行版，在技术内容上向 UN R48-09 进一步靠拢，整体技术对齐度从约 70\% 提升至约 90\%。但仍保留部分中国特色条款（如特定车型的特殊规定、地方差异化要求等）。

\textbf{四、合规影响：} 草案版一旦正式发布，整车厂和灯具供应商需要：
\begin{itemize}[leftmargin=2em]
    \item 重新评估现有产品的合规性；
    \item 更新测试验证流程；
    \item 调整设计参数（如光度、色度、安装尺寸）；
    \item 重新申请型式认证或扩展认证。
\end{itemize}

\section*{主要建议}

针对整车厂、灯具供应商和检测认证机构，本报告提出以下建议：
\begin{enumerate}[leftmargin=2em]
    \item \textbf{提前规划}：密切关注草案版本的最终发布，及时调整产品规划；
    \item \textbf{技术预研}：对 ADB、矩阵式 LED 等新型照明技术进行技术储备；
    \item \textbf{认证协同}：与认证机构保持密切沟通，预判认证流程的变化；
    \item \textbf{供应链协调}：与上游 LED 芯片、控制器供应商协同应对新要求；
    \item \textbf{国际化布局}：同步关注 UN R48 后续系列修订，做好国际认证对接。
\end{enumerate}

\section*{报告结构}

本报告共分 13 章正文及 4 个附录：
\begin{itemize}[leftmargin=2em]
    \item 第 1--2 章：法规背景与适用范围；
    \item 第 3--4 章：总体架构与术语对比；
    \item 第 5--8 章：一般要求与具体装置对比；
    \item 第 9--10 章：光度色度与测试方法对比；
    \item 第 11--13 章：变化点总结、合规影响与结论建议；
    \item 附录 A--D：对照表、装置代码、关键数据差异、术语对照。
\end{itemize}

\vspace{2cm}
\noindent\rule{\textwidth}{0.5pt}

\textbf{关键词：} GB 4785；UN R48；汽车外部照明；光信号装置；安装规定；自适应前照灯；自适应远光灯；昼间行驶灯；型式认证；横向对比

\clearpage